\subsection{Compact operators and spectral decomposition}\label{subsec:spectral-preliminaries-operators}
\paraid{p-17e565d7cc93}
For a Hilbert space
$\HilbertDomain $,
we call a bounded linear operator
$\CompactOperator \colon \HilbertDomain \to \HilbertDomain $
\emph{compact} if the image of the closed unit ball of
$\HilbertDomain$
has compact closure.
We call
$\CompactOperator$
\emph{self-adjoint} if for every
$\HilbertFunction,\HilbertTestFunction\in\HilbertDomain$
we have
\[
 \langle \CompactOperator \HilbertFunction ,\HilbertTestFunction \rangle_\HilbertDomain =\langle \HilbertFunction ,\CompactOperator \HilbertTestFunction \rangle_\HilbertDomain.
\]
A nonzero vector
$\HilbertFunction $
is an \emph{eigenvector} for the eigenvalue
$\Eigenvalue $
if
$\CompactOperator \HilbertFunction =\Eigenvalue  \HilbertFunction $.
For an operator on a function space we also call it an eigenfunction.

\paraid{p-e197c5fac09a}
For a real Hilbert space
$\HilbertDomain$,
write
$\HilbertDomain_{\ComplexNumbers}=\HilbertDomain\oplus\ImaginaryUnit\HilbertDomain$
for its complexification.
We denote the complex-linear extension of
$\CompactOperator$
by the same symbol.
For
$\HilbertFunction,\HilbertTestFunction\in\HilbertDomain$,
set
\[
 \CompactOperator(\HilbertFunction+\ImaginaryUnit\HilbertTestFunction)
 =\CompactOperator\HilbertFunction+\ImaginaryUnit\CompactOperator\HilbertTestFunction.
\]
For a complex Banach space,
write
$\IdentityMap$
for its identity operator.
For a bounded operator
$\CompactOperator$
on that space,
its \emph{spectrum}
$\spec(\CompactOperator)$
is the set of complex numbers
$\Eigenvalue$
for which
$\Eigenvalue\IdentityMap-\CompactOperator$
has no bounded inverse.
For an operator on a real Hilbert space,
we use the spectrum of its
complex-linear extension to
$\HilbertDomain_{\ComplexNumbers}$.
For a self-adjoint operator this set is real.
In finite dimension it is exactly the set of eigenvalues.
In infinite dimension,
the number  $0$ can belong to the spectrum without being an
eigenvalue.

\begin{theorem}[{Riesz--Schauder and the self-adjoint spectral theorem, \cite[Theorems~14.19 and~15.22]{Muscat2024}}]\label{prop:spectral-basics}
\paraid{p-0edefea64bb3}
Let
$\CompactOperator $
be a compact self-adjoint operator on a separable real Hilbert space
$\HilbertDomain $.
Then every nonzero spectral value is a real eigenvalue with finite-dimensional
eigenspace.
For each
$\SpectralCutoff >0$,
there are only finitely many eigenvalues,
counted with multiplicity,
with absolute value greater than
$\SpectralCutoff $.
There is a finite or countable orthonormal family
$\{\Eigenfunction _\CoordinateIndex \}_{\CoordinateIndex \in \EigenIndexSet }$
indexed by a subset
$\EigenIndexSet\subset\NonnegativeIntegers$,
with nonzero eigenvalues
$\{\Eigenvalue _\CoordinateIndex \}_{\CoordinateIndex \in \EigenIndexSet }$
such that
\[
 \HilbertDomain =\ker \CompactOperator \ \mathbin{\oplus}\
       \overline{\LinearSpan }\{\Eigenfunction _\CoordinateIndex \mid \CoordinateIndex \in \EigenIndexSet \}.
\]
Every
$\HilbertFunction \in \HilbertDomain $
admits the decomposition
\[
 \HilbertFunction =\HilbertFunction _0+\sum_{\CoordinateIndex \in \EigenIndexSet }\langle \HilbertFunction ,\Eigenfunction _\CoordinateIndex \rangle \Eigenfunction _\CoordinateIndex .
\]
The operator acts on this decomposition by
\[
 \CompactOperator \HilbertFunction =\sum_{\CoordinateIndex \in \EigenIndexSet }\Eigenvalue _\CoordinateIndex \langle \HilbertFunction ,\Eigenfunction _\CoordinateIndex \rangle \Eigenfunction _\CoordinateIndex .
\]
The remaining component satisfies
\[
 \HilbertFunction _0\in\ker \CompactOperator .
\]
Both sums converge in
$\HilbertDomain $.
If
$\EigenIndexSet $
is infinite,
its eigenvalues tend to
$0$
in any enumeration.
\end{theorem}

\paraid{p-c81130e9d92f}
This is the self-adjoint case of the spectral theorem for compact normal
operators
\cite[Theorems~14.19 and~15.22]{Muscat2024}.
In particular,
a nonzero compact self-adjoint operator has a nonzero eigenvalue.
